\documentclass[11pt]{article}
\usepackage[margin=1in]{geometry}
\usepackage[T1]{fontenc}
\usepackage{lmodern,amsmath,amssymb,amsthm,mathtools,hyperref}
\setlength{\emergencystretch}{3em}
\newtheorem{theorem}{Theorem}[section]
\newtheorem{lemma}[theorem]{Lemma}
\newtheorem{corollary}[theorem]{Corollary}
\newtheorem{proposition}[theorem]{Proposition}
\newcommand{\E}{\mathbb E}
\newcommand{\tr}{\operatorname{tr}}
\title{Sharp cubic encodings of known Hadamard matrices\\on minimal binary seeds}
\author{Joshua Mark\thanks{ORCID: \href{https://orcid.org/0009-0003-5905-2382}{0009-0003-5905-2382}. AI tools assisted with derivation, independent review, computation and exposition.}}
\date{20 November 2026 --- research draft; priority not established}
\begin{document}\maketitle
\section{Minimal binary seeds require substantial nonlinear observables}
\label{sec:nonlinear-seeds}
We distinguish Boolean algebraic degree over $\mathbb F_2$ from real
Walsh degree. A sign observable is written $(-1)^{f(z)}$, and its degree
means the degree of the unique squarefree Boolean polynomial $f$.
The constructions below encode the classical Paley matrices; they do
not give new Hadamard orders. Paley's primary paper is
\emph{On Orthogonal Matrices} (1933),
\url{https://doi.org/10.1002/sapm1933121311}. The standard binary-code
viewpoint also appears in Orrick, \emph{Switching operations for Hadamard
matrices}, \url{https://arxiv.org/abs/math/0507515}.
All encoding bounds below are proved explicitly; novelty is unestablished.

\begin{proposition}[Affine observable budget]
\label{prop:nonlinear-affine-budget}
Suppose an order-$N$ normalized Hadamard matrix is encoded by distinct
labels $S\subset\mathbb F_2^r$, with uniform weight $1/N$ on each label.
If $N\equiv4\pmod8$ and $r<N-2$, the number $k$ of nonconstant columns
that extend to affine Boolean functions satisfies
\[
 k\le\min\{r,\,2^r-N-1\}\qquad(N<2^r).
\]
Thus a minimal four-bit encoding of $H_{12}$ needs at least eight
nonaffine observables, and a minimal five-bit encoding of $H_{20}$
needs at least fourteen.
\end{proposition}
\begin{proof}
Write $B=(J-H_{\rm core})/2$. Over $\mathbb F_2$, the core rows have
even weight $N/2$ and distinct rows intersect in odd weight $N/4$.
Therefore $BB^T=J+I$, of rank $N-2$. Also $B\mathbf1=0$, so
$\operatorname{rank}B=N-2$ and its only column relation is the sum
of all $N-1$ columns. Any proper subset is independent. Affine functions
vanishing at the label of the normalized first row form a space of
dimension $r$. Hence at most $r$ columns can be affine.

For the other bound, affine sign functions giving distinct nonconstant
orthogonal columns have distinct nonzero linear parts. Together with
one, they are orthogonal on the whole binary cube. Their inner products
vanish on $S$, so also on its complement $T$. Thus $k+1$ nonzero
orthogonal real functions live on $|T|=2^r-N$ points, giving the bound.
\end{proof}

Let $\mathrm{RM}(d,r)$ denote the space of Boolean polynomials of degree
at most $d$ on $r$ bits; its dimension is $\sum_{j=0}^d\binom rj$.
\begin{lemma}[Complement criterion for low-degree extensions]
\label{lem:nonlinear-extension}
Every Boolean function on $S\subset\mathbb F_2^r$ has a degree-at-most-$d$
extension if and only if there is no nonzero polynomial of degree at
most $r-d-1$ supported entirely on $S$. In particular, cubic extensions
exist on any proper subset of the four-bit cube. On five bits they
exist whenever the complement of $S$ affinely spans the cube.
\end{lemma}
\begin{proof}
The binary inner product of monomials of total combined degree less
than $r$ is zero: summing over any missing variable gives zero modulo
two. Dimension counting then gives
$\mathrm{RM}(d,r)^\perp=\mathrm{RM}(r-d-1,r)$.
A nonzero left-kernel vector of restriction to $S$, extended by zero
off $S$, is exactly such a dual polynomial. The particular cases use
that a nonzero constant cannot vanish on a nonempty complement, and
that an affine function vanishing on an affine spanning set is zero.
\end{proof}

\begin{theorem}[Sharp cubic benchmarks at orders twelve and twenty]
\label{thm:nonlinear-cubic}
For $N=12$ and $N=20$, the least possible maximum Boolean degree of
all $N-1$ observables, among encodings on the minimal number
$r=\lceil\log_2N\rceil$ of binary seed bits, is exactly three.
The affine budgets above are simultaneously sharp: the examples have
exactly three affine observables at order twelve and five at order twenty.
\end{theorem}
\begin{proof}
For $N\equiv4\pmod8$, the normalized binary columns together with the
constant function span a space of dimension $N-1$: the binary columns
have rank $N-2$ and vanish at the first row. Thus any common degree
bound $d$ requires
$\sum_{j=0}^d\binom rj\ge N-1$.
At $N=20,r=5$, quadratics have dimension $16<19$, proving the lower
bound three.

At $N=12,r=4$, the quadratic dimension is eleven, exactly the required
rank. Suppose all observables were quadratic. Restriction of
$\mathrm{RM}(2,4)$ to the twelve labels would have rank eleven and equal
the span of the binary Hadamard columns and the constant function.
Every vector in the latter span has even weight: the column weights
are six and the constant has weight twelve. Hence all coordinate
functions have even weight on $S$. Their full-cube weights are even,
so the four omitted labels have xor zero. Four distinct labels with xor
zero form an affine plane. Its indicator is a nonzero quadratic
polynomial vanishing on $S$, contradicting the restriction rank eleven.
This rules out quadratics for every labeling, without exhaustive search.

For explicit upper bounds, use Paley's normalized order-$p+1$ matrix
for $p=11,19$:
\[
 H_{\infty,j}=H_{i,\infty}=1,\qquad
 H_{i,j}=\chi(j-i)-\mathbf1_{i=j},
\]
where $\chi$ is the Legendre symbol with $\chi(0)=0$.
The first three nonconstant columns of $H_{12}$ give eight binary
patterns, four occurring twice and four once. Label rows by these
three bits and a fourth discriminator, choosing zero for the first
occurrence and one for the second. In row order $(\infty,0,\ldots,10)$
the resulting four-bit integer labels are
\[
 0,5,3,6,13,11,7,14,4,8,1,2.
\]
The first three observables are affine. Every other observable has a
cubic extension by the lemma. The affine budget proves three is maximal.

The first five nonconstant columns of $H_{20}$ separate all twenty
rows, with five-bit labels
\[
 0,13,27,22,12,25,19,7,15,30,29,26,21,10,20,8,16,1,3,6.
\]
The complement contains the affine basis
$2,4,5,9,11,17$: their differences from $2$ have binary rank five.
The lemma gives cubic extensions of every observable, retaining the
first five as coordinate functions. The affine budget proves optimality.
All displayed finite certificates and extension polynomials are checked
by original exact binary elimination in the accompanying script.
\end{proof}

\paragraph{A reusable nonlinear seed rule and its limit.}
For any odd prime power $q\equiv3\pmod4$, Paley's quadratic-character
formula over $\mathbb F_q$ supplies an extremal $(q+1)$-atom distribution.
Inject its row labels into a binary cube, use the uniform measure on
those labels, and solve the Boolean restriction equations for the
observables. Lemma~\ref{lem:nonlinear-extension} precisely certifies when
a chosen degree budget can express all outputs. The order-twelve and
order-twenty examples show that cubic outputs genuinely cross the affine
seed obstruction at minimal seed length; they are exact extremal laws,
not approximate moment fits. They still inherit Paley's restriction on
$q$, so this rule supplies no distribution at a currently unknown order.

\begin{corollary}[Fixed algebraic degree cannot suffice asymptotically]
\label{cor:nonlinear-growing-degree}
For any Hadamard order $N\ge12$ with $N\equiv4\pmod8$, and minimal
binary seed length $r=\lceil\log_2N\rceil$, every encoding has maximum
Boolean degree at least $\lceil r/2\rceil$. Conversely, every injection
of its rows into the $r$-bit cube admits extensions of degree at most
$r-1$. Thus a fixed-degree minimal-seed model cannot represent extremal
laws in arbitrarily large such orders, even where Hadamard matrices
are already known.
\end{corollary}
\begin{proof}
The rank bound requires $\sum_{j=0}^d\binom rj\ge N-1$.
Here $N-1>2^{r-1}$, while binomial symmetry gives
$\sum_{j=0}^{\lfloor(r-1)/2\rfloor}\binom rj\le2^{r-1}$.
This proves the lower bound. The proper subset $S$ has nonempty
complement, and Lemma~\ref{lem:nonlinear-extension} with $d=r-1$
gives the upper bound.
\end{proof}

\end{document}
